\section{Population ReLU Kernel on the Circle}

Throughout, $(\cdot|\cdot)$ denotes the Euclidean inner product on $\mathbb R^d$.

\subsection{Population ReLU Feature Kernel}

Let $\sigma(t)=\max\{0,t\}$ denote the ReLU activation and let
$w\sim\mathcal N(0,I_d)$.
We define the \emph{population ReLU feature kernel} by
\begin{equation}
	\label{eq:population-kernel-def}
	K(x,y)
	\;:=\;
	\mathbb E_w\big[\sigma((w|x))\sigma((w|y))\big],
	\qquad x,y\in\mathbb R^d.
\end{equation}
This kernel is the covariance of the random feature map
$\phi_w(x)=\sigma((w|x))$ and is therefore positive semidefinite.

Let $\gamma_d$ denote the standard Gaussian measure on $\mathbb R^d$,
\[
	d\gamma_d(w)
	=
	(2\pi)^{-d/2}\exp\!\Big(-\frac{\|w\|^2}{2}\Big)\,dw.
\]
Then \eqref{eq:population-kernel-def} can be written as the Gaussian integral
\begin{equation}
	\label{eq:population-kernel-integral}
	K(x,y)
	=
	\int_{\mathbb R^d}
	\sigma((w|x))\sigma((w|y))\,
	d\gamma_d(w).
\end{equation}

Using $\sigma(t)=t\,\mathbf 1_{\{t>0\}}$, this becomes
\begin{equation}
	\label{eq:relu-indicator-integral}
	K(x,y)
	=
	\int_{\mathbb R^d}
	(w|x)(w|y)\,
	\mathbf 1_{\{(w|x)>0\}}\,
	\mathbf 1_{\{(w|y)>0\}}\,
	d\gamma_d(w).
\end{equation}

\subsection{Reduction to a Two-Dimensional Gaussian}

Let $V=\mathrm{span}\{x,y\}\subset\mathbb R^d$.

\begin{lemma}[Orthogonal decomposition]
	\label{lem:orth-decomp}
	For every $w\in\mathbb R^d$ there exist unique vectors
	$w_0\in V$ and $w_\perp\in V^\perp$ such that
	\[
		w = w_0 + w_\perp,
		\qquad
		(w_0|w_\perp)=0.
	\]
\end{lemma}

\begin{proof}
	Choose an orthonormal basis $\{e_1,\dots,e_k\}$ of $V$ and extend it to an
	orthonormal basis $\{e_1,\dots,e_d\}$ of $\mathbb R^d$.
	Writing $w=\sum_{i=1}^d (w|e_i)e_i$, define
	\[
		w_0 := \sum_{i=1}^k (w|e_i)e_i,
		\qquad
		w_\perp := \sum_{i=k+1}^d (w|e_i)e_i.
	\]
	Then $w=w_0+w_\perp$, $(w_0|w_\perp)=0$, and uniqueness follows from the
	uniqueness of the orthonormal expansion.
\end{proof}

Since $x,y\in V$, we have $(w_\perp|x)=(w_\perp|y)=0$, hence
\[
	(w|x)=(w_0|x),
	\qquad
	(w|y)=(w_0|y).
\]
Thus the integrand in \eqref{eq:relu-indicator-integral} depends only on $w_0$.

\begin{lemma}[Gaussian factorization]
	\label{lem:gaussian-factor}
	If $w\sim\mathcal N(0,I_d)$, then $w_0$ and $w_\perp$ are independent Gaussian
	vectors with
	\[
		w_0\sim\mathcal N(0,I_V),
		\qquad
		w_\perp\sim\mathcal N(0,I_{V^\perp}).
	\]
\end{lemma}

\begin{proof}
	In an orthonormal basis adapted to $V\oplus V^\perp$, the coordinates of $w$
	are independent standard Gaussians. Grouping coordinates yields the claim.
\end{proof}

Hence integration over $V^\perp$ contributes a factor of one, and the kernel
reduces to a Gaussian integral over $V$.

\subsection{Choice of Coordinates and Angle Representation}

Assume $x\neq 0$ and define
\[
	u := \frac{x}{\|x\|}.
\]
If $y$ is not collinear with $x$, define
\[
	v := \frac{y-(y|u)u}{\|y-(y|u)u\|}.
\]
Then $\{u,v\}$ is an orthonormal basis of $V$.

Define the angle $\theta\in[0,\pi]$ by
\[
	\cos\theta := \frac{(x|y)}{\|x\|\|y\|}.
\]

\begin{lemma}[Angle decomposition]
	\label{lem:angle-decomp}
	With $u,v$ as above,
	\[
		x=\|x\|u,
		\qquad
		y=\|y\|(\cos\theta\,u+\sin\theta\,v).
	\]
\end{lemma}

\begin{proof}
	Since $y\in V$,
	\[
		y=(y|u)u+(y|v)v.
	\]
	Now $(y|u)=\|y\|\cos\theta$. By Pythagoras,
	\[
		(y|v)^2=\|y\|^2-(y|u)^2=\|y\|^2\sin^2\theta,
	\]
	and by construction $(y|v)\ge 0$, hence $(y|v)=\|y\|\sin\theta$.
\end{proof}

\subsection{Gaussian Coordinates in the Plane}

\begin{lemma}[Coordinates of $w_0$]
	\label{lem:w0-coords}
	There exist unique scalars $Z_1,Z_2$ such that
	\[
		w_0=Z_1u+Z_2v,
		\qquad
		Z_1=(w|u),\; Z_2=(w|v).
	\]
\end{lemma}

\begin{proof}
	Since $\{u,v\}$ is an orthonormal basis of $V$, write $w_0=Z_1u+Z_2v$.
	Taking inner products with $u$ and $v$ gives $Z_1=(w_0|u)$ and $Z_2=(w_0|v)$.
	Because $w=w_0+w_\perp$ and $u,v\perp w_\perp$, these equal $(w|u)$ and $(w|v)$.
\end{proof}

\begin{lemma}[Distribution of $(Z_1,Z_2)$]
	\label{lem:z-gaussian}
	Let $w\sim\mathcal N(0,I_d)$. Then
	\[
		(Z_1,Z_2)\sim\mathcal N(0,I_2),
	\]
	and in particular $Z_1,Z_2$ are independent standard normal random variables.
\end{lemma}

\begin{proof}
	For any $a\in\mathbb R^d$, $(a|w)\sim\mathcal N(0,\|a\|^2)$.
	Thus $Z_1,Z_2\sim\mathcal N(0,1)$.
	Moreover,
	\[
		\mathrm{Cov}(Z_1,Z_2)=\mathbb E[(w|u)(w|v)]=(u|v)=0.
	\]
	Since any linear combination $aZ_1+bZ_2=(w|au+bv)$ is Gaussian, the vector
	$(Z_1,Z_2)$ is jointly Gaussian. Zero covariance implies independence.
\end{proof}

Combining the above,
\[
	(w|x)=\|x\|Z_1,
	\qquad
	(w|y)=\|y\|(\cos\theta\,Z_1+\sin\theta\,Z_2).
\]

\subsection{Evaluation of the Gaussian Integral}

Substituting into \eqref{eq:relu-indicator-integral} yields
\begin{equation}
	\label{eq:kernel-2d}
	K(x,y)
	=
	\|x\|\|y\|\,
	\mathbb E\Big[
		Z_1(\cos\theta\,Z_1+\sin\theta\,Z_2)
		\mathbf 1_{\{Z_1>0\}}
		\mathbf 1_{\{\cos\theta Z_1+\sin\theta Z_2>0\}}
		\Big],
\end{equation}
where $(Z_1,Z_2)\sim\mathcal N(0,I_2)$.

The joint density of $(Z_1,Z_2)$ is
\[
	\frac{1}{2\pi}\exp\!\Big(-\frac{Z_1^2+Z_2^2}{2}\Big).
\]
Introduce polar coordinates $(Z_1,Z_2)=(R\cos\varphi,R\sin\varphi)$ with
$R\ge 0$ and $\varphi\in[0,2\pi)$. Then the Gaussian measure becomes
\[
	\frac{1}{2\pi}e^{-R^2/2}R\,dR\,d\varphi,
\]
and
\[
	\cos\theta\,Z_1+\sin\theta\,Z_2 = R\cos(\varphi-\theta).
\]

The indicator constraints are equivalent to
\[
	\cos\varphi>0,
	\qquad
	\cos(\varphi-\theta)>0,
\]
which define the angular interval
\[
	\varphi\in(0,\pi)\cap(\theta,\theta+\pi),
\]
of length $\pi-\theta$.

Separating radial and angular variables gives
\[
	\mathbb E[\cdots]
	=
	\frac{1}{2\pi}
	\int_0^\infty R^3 e^{-R^2/2}\,dR
	\int_{\mathcal I_\theta}
	\cos\varphi\cos(\varphi-\theta)\,d\varphi,
\]
where $\mathcal I_\theta$ denotes the above interval.
The radial integral evaluates to
\[
	\int_0^\infty R^3 e^{-R^2/2}\,dR = 2,
\]
and the angular integral yields
\[
	\int_{\mathcal I_\theta}\cos\varphi\cos(\varphi-\theta)\,d\varphi
	=
	\frac{1}{2}\big(\sin\theta+(\pi-\theta)\cos\theta\big).
\]

\begin{theorem}[Arc-cosine kernel]
	\label{thm:arc-cosine}
	For all $x,y\in\mathbb R^d$,
	\[
		K(x,y)
		=
		\frac{\|x\|\|y\|}{2\pi}
		\big(\sin\theta+(\pi-\theta)\cos\theta\big),
		\qquad
		\theta=\arccos\!\Big(\frac{(x|y)}{\|x\|\|y\|}\Big).
	\]
\end{theorem}

\subsection{Restriction to the Unit Circle}

Let $x(\phi)=(\cos\phi,\sin\phi)\in S^1$.
Then $(x(\phi)|x(\psi))=\cos(\phi-\psi)$ and
\[
	K(x(\phi),x(\psi))=k(\theta(\phi-\psi)),
	\qquad
	k(\theta)=\frac{1}{2\pi}\big(\sin\theta+(\pi-\theta)\cos\theta\big).
\]

\subsection{Population Integral Operator}

Let $\mu$ be the uniform probability measure on $[0,2\pi)$.
Define
\[
	(Tf)(\phi)
	=
	\int_0^{2\pi}K(x(\phi),x(\psi))f(\psi)\,d\mu(\psi).
\]
Then $T$ is a circular convolution operator.

\subsection{Discretization of the Population}

Let $\phi_1,\dots,\phi_n\stackrel{iid}{\sim}\mathrm{Unif}[0,2\pi)$ and define
\[
	(K_n)_{ij}=K(x(\phi_i),x(\phi_j)).
\]
Equivalently,
\[
	(T_nf)_i=\frac1n\sum_{j=1}^nK(x(\phi_i),x(\phi_j))f(\phi_j).
\]
This is a Monte Carlo discretization of the population operator $T$,
isolating finite-sample effects while remaining in the infinite-population
(kernel) regime.
